\section{A complete two-descent and the actual square--cube exclusion}
\label{zd2x-section}

This section proves the rank-zero hypothesis left explicit in
Lemma~\ref{z2x-rank-zero-points}. It uses a descent on the entire
Jacobian, including closed points over finite local extensions.
No software rank estimate, analytic rank, or unproved saturation
claim enters the argument.

\begin{lemma}[The monic model and its actual descent classes]
\label{zd2x-model-descent}
The change \(X=-3s,\ Y=9y\) identifies the smooth projective curve
of~\eqref{z2x-quintic} with
\begin{equation}
 C:Y^2=F(X),\qquad
 F(X)=X(X-3)g(X),\qquad g(X)=X^3-27X+27.
 \label{zd2x-monic}
\end{equation}
Let \(O\) be its unique point at infinity, \(J=\operatorname{Jac}(C)\),
and \(K=\mathbb Q(\theta)\), where \(\theta^3-3\theta-1=0\).
Writing \(\alpha=-3\theta\), its descent algebra is
\begin{equation}
 L=\mathbb Q[T]/(F(T))=\mathbb Q\times\mathbb Q\times K,\qquad
 T\longmapsto(0,3,\alpha),
 \label{zd2x-algebra}
\end{equation}
with norm \(N(q_0,q_3,\beta)=q_0q_3N_{K/\mathbb Q}(\beta)\).
The rational two-torsion classes
\(D_0=[(0,0)-O]\), \(D_3=[(3,0)-O]\) have descent images
\begin{equation}
 \delta(D_0)=(-1,-3,3\theta),\qquad
 \delta(D_3)=(3,-1,3(\theta+1))
 \quad\bmod L^{*2}.
 \label{zd2x-torsion-images}
\end{equation}
\end{lemma}
\begin{proof}
Direct substitution gives \(F(-3s)=81[-3s(s+1)(s^3-3s-1)]\).
The rational affine change and its inverse extend over the unique
smooth points at infinity. For clarity, the standard odd-degree
descent input is used in the following precise form, for every
characteristic-zero field \(k\):
\begin{equation}
 \delta_k:J(k)\longrightarrow
 \ker\bigl(N:L_k^*/L_k^{*2}\longrightarrow k^*/k^{*2}\bigr),
 \qquad \ker\delta_k=2J(k).
 \label{zd2x-delta-kernel}
\end{equation}
The map is natural in \(k\). On a degree-zero rational divisor
avoiding the Weierstrass points and \(O\), it evaluates \(X-T\),
using closed-point norms and the actual divisor multiplicities.
Every class of \(J(k)\) admits such a representative by the
rational base point and divisor moving. Thus this is a map on
all of \(J(k)\), not merely the image of \(C(k)\).
These are the construction preceding Lemma~4.1 and Lemma~4.1 of
\cite{StollTwoDescent2001}. The monic odd-degree hypotheses
apply to~\eqref{zd2x-monic}; no fake-descent kernel for an
even-degree model is being used.

For a nonbranch point \(P\), the value on \(P-O\) is \(X(P)-T\).
At a rational branch point \((a,0)\), its component at \(T=a\)
is \(F'(a)\), and the other components are \(a-T\), modulo
squares; this is \cite[Lemma~4.3]{StollTwoDescent2001}.
Here
\[
 g(0)=27,\qquad g(3)=-27,\qquad F'(0)=F'(3)=-81.
\]
Substitution gives~\eqref{zd2x-torsion-images}; in particular
the first component at \(D_0\) simplifies from \(-81\) to
\(-1\), not to \(1\).
\end{proof}

\begin{lemma}[The cubic field and all positive unit square classes]
\label{zd2x-field}
One has
\[
 \mathcal O_K=\mathbb Z[\theta],\qquad
 \operatorname{Disc}(K)=81,\qquad
 \operatorname{Cl}(\mathcal O_K)=1.
\]
The element \(\pi=\theta-1\), of norm three, generates the unique
prime above three, whose residue degree is one. Every totally
positive unit of \(\mathcal O_K\) is a square.
\end{lemma}
\begin{proof}
The polynomial \(h(U)=U^3-3U-1\) has no rational root, since its
possible integral roots \(1,-1\) both fail. It has exactly three
real roots, one in each interval
\begin{equation}
 1<\theta_1<2,\qquad -2<\theta_2<-1,\qquad -1<\theta_3<0.
 \label{zd2x-root-intervals}
\end{equation}
The intermediate value theorem supplies all three, and its
discriminant is \(81\). If \(\mathbb Z[\theta]\) had nontrivial
index in \(\mathcal O_K\), the index-discriminant identity
would force \(\operatorname{Disc}(K)\le9\): the square of that
index divides \(81\). The Minkowski bound for a totally real
cubic field is \((2/9)\sqrt{\operatorname{Disc}(K)}\), by
\cite[Theorem~4.3]{MilneANT2020}. It would then be at most
\(2/3\), contradicting the positive integral norm of an ideal
representative required by that bound. Thus the index is one.

The actual bound is now exactly two. Modulo two, \(h\) is
\(U^3+U+1\), which is irreducible since it has no root in
\(\mathbb F_2\). In the established maximal monogenic order,
two is therefore inert and there is no ideal of norm two.
Every ideal class has an integral representative of norm one
and is trivial.

The shifted polynomial
\[
 h(1+U)=U^3+3U^2-3
\]
is Eisenstein at three. Thus there is a unique prime above
three, of ramification degree three and residue degree one.
The element \(\pi\) has norm three and generates that prime.
This also proves that \(g\) is irreducible over \(\mathbb Q_3\).

The units \(-1,\theta,\theta+1\) have norms \(-1,1,-1\).
Their signatures in the ordering~\eqref{zd2x-root-intervals}
are
\begin{equation}
 \begin{array}{c|ccc}
  &\theta_1&\theta_2&\theta_3\\ \hline
  -1&-&-&-\\
  \theta&+&-&-\\
  \theta+1&+&-&+
 \end{array}
 \label{zd2x-unit-signatures}
\end{equation}
and are independent over \(\mathbb F_2\).
Dirichlet's theorem \cite[Theorem~5.1]{MilneANT2020} gives
\(\mathcal O_K^*\simeq\{\pm1\}\times\mathbb Z^2\).
Hence its quotient by squares has dimension three, and the
signature map onto \(\{\pm1\}^3\) is an isomorphism. A totally
positive unit is consequently a square. This argument needs
only the displayed classes modulo squares; it makes no claim
that \(\theta,\theta+1\) form a fundamental unit basis over
\(\mathbb Z\), and invokes no regulator or class-group oracle.
\end{proof}

\begin{lemma}[All closed-point evaluations have even valuations away from three]
\label{zd2x-even-valuations}
For every \(Q\in J(\mathbb Q)\), each component of \(\delta(Q)\)
has even ideal valuation at every prime not above three.
Consequently its square class has a representative
\begin{equation}
 (q_0,q_3,u\pi^e),\qquad
 q_0,q_3\in\{1,-1,3,-3\},\quad
 e\in\{0,1\},\quad u\in\mathcal O_K^* .
 \label{zd2x-global-classes}
\end{equation}
The valuation statement includes the prime two.
\end{lemma}
\begin{proof}
We prove the stronger statement on all of \(J(\mathbb Q_\ell)\)
for each prime \(\ell\ne3\). The discriminants of \(X(X-3)\)
and \(g\) are \(3^2,3^{10}\), and their resultant is
\(g(0)g(3)=-3^6\). Thus
\[
 \operatorname{Disc}(F)=3^{24}.
\]
It follows that all five roots of \(F\) are integral and have
unit pairwise differences in some finite unramified extension
\(U/\mathbb Q_\ell\). This remains true when \(\ell=2\).

Let \(P\) be any closed point over \(\mathbb Q_\ell\), away
from the excluded support of~\eqref{zd2x-delta-kernel}, with
finite residue field \(E/\mathbb Q_\ell\). Write \(x=X(P)\),
\(y=Y(P)\). The compositum \(E'=EU\) is unramified over \(E\).
In its normalized integer-valued valuation, if \(v(x)<0\),
then \(v(x-r_i)=v(x)\) for all five roots and
\[
 2v(y)=5v(x).
\]
Hence every factor valuation is even. If \(v(x)\ge0\), all
five factors are integral and at most one can be a nonunit,
since the roots have unit differences. Its valuation is the
whole valuation of \(F(x)=y^2\), and is again even.

We spell out the norm transfer to avoid a restriction to
rational points. Each field factor \(M\) of \(L_{\mathbb Q_\ell}\)
is unramified over \(\mathbb Q_\ell\). Write
\(E\otimes_{\mathbb Q_\ell}M=\prod_j V_j\).
Each \(V_j/E\) is unramified. Extending it by \(U\), still
unramified, splits all roots and gives the preceding valuation
test. Thus \(x-T\) in each \(V_j\) has even normalized valuation.
Its norm to \(M\) multiplies that valuation by the residue
degree \(f(V_j/M)\), preserving evenness. The \(M\)-component
of the closed-point evaluation is the product of these norms.
Positive or negative divisor multiplicities preserve parity.
Using a moved divisor for any class of \(J(\mathbb Q_\ell)\)
proves the local assertion on the whole Jacobian. Naturality
of~\eqref{zd2x-delta-kernel} proves the global assertion.

Only even ideal valuation is claimed here. In particular it
does not say that a quadratic square-class extension is
unramified at two, nor assume good reduction of the equation
at two. This direct argument supplies the condition at that
prime instead of deleting it from a generic Selmer-prime list.

The rational factors with even valuations away from three
give the four choices for each \(q_i\). In \(K\), the
fractional principal ideal of a representative \(\beta\)
has the form
\[
 (\beta)=(\pi)^e\mathfrak a^2,\qquad e\in\{0,1\}.
\]
An even exponent at \(\pi\), including a negative one, is
absorbed into \(\mathfrak a\). Lemma~\ref{zd2x-field} gives
\(\mathfrak a=(\gamma)\), so
\(\beta=u\pi^e\gamma^2\) for a unit \(u\).
This proves~\eqref{zd2x-global-classes} with all denominators
included.
\end{proof}

\begin{lemma}[The full three-adic image and global matching]
\label{zd2x-three-adic}
The four elements of \(\delta_{\mathbb Q_3}(J(\mathbb Q_3))\)
are the images of \(0,D_0,D_3,D_0+D_3\).
Multiplying any global image by a suitable one of these
rational torsion images leaves an actual global image of
the form \((1,1,u)\), with \(u\in\mathcal O_K^*\) and
\(N_{K/\mathbb Q}(u)=1\).
\end{lemma}
\begin{proof}
There are exactly three irreducible factors of \(F\) over
\(\mathbb Q_3\), because \(g\) is irreducible there.
The odd-degree branch-divisor description
\cite[Lemma~4.3(3)]{StollTwoDescent2001} gives
\[
 \dim_{\mathbb F_2}J(\mathbb Q_3)[2]=3-1=2.
\]
The quotient by doubling has the same order, by
\cite[Lemma~4.4(1)]{StollTwoDescent2001}. One may also see
this directly. Take a compact open formal subgroup \(U\)
of \(J(\mathbb Q_3)\) on which doubling is bijective: its
linear term is the three-adic unit \(2\), and a sufficiently
small formal chart gives its inverse by contraction.
The quotient \(A=J(\mathbb Q_3)/U\) is finite. The groups
\(J(\mathbb Q_3)/2J(\mathbb Q_3)\) and \(J(\mathbb Q_3)[2]\)
identify respectively with \(A/2A\) and \(A[2]\).
For the second assertion a lift of an element of \(A[2]\)
is uniquely corrected by a half in \(U\).
The equality \(|A/2A|=|A[2]|\) proves the order four.

The first components \(-1,3\) of
\eqref{zd2x-torsion-images} are independent in
\(\mathbb Q_3^*/\mathbb Q_3^{*2}\): the first is a nonsquare
unit and the second has odd valuation. Thus the two local
images are independent and exhaust this four-element group.
Furthermore, the four rational square classes
\(\{1,-1,3,-3\}\) map injectively into that local square-class
group. By Lemma~\ref{zd2x-even-valuations} and naturality,
the first two global components are therefore one of
\begin{equation}
 (1,1),\qquad(-1,-3),\qquad(3,-1),\qquad(-3,3)
 \quad\bmod\mathbb Q^{*2}.
 \label{zd2x-local-matching}
\end{equation}
Multiplication by the corresponding actual rational torsion
image leaves the first two components globally square.

For this new global image write the remaining component as
\(u\pi^e\) modulo squares, as in~\eqref{zd2x-global-classes}.
The norm condition in~\eqref{zd2x-delta-kernel} says
\(N(u)3^e\) is a rational square. Since \(N(u)=\pm1\)
and \(e=0\) or \(1\), it follows that \(e=0\), \(N(u)=1\).
The residue degree one of the unique prime above three
is explicit in the equality \(N(\pi)=3\).
\end{proof}

\begin{theorem}[Rank zero and the complete three-point quotient]
\label{zd2x-rank-and-points}
For the curve~\eqref{z2x-quintic},
\[
 J(\mathbb Q)\simeq(\mathbb Z/2\mathbb Z)^2,\qquad
 H(\mathbb Q)=\{P_0,P_{-1},P_\infty\}.
\]
\end{theorem}
\begin{proof}
We first kill the remaining unit in
Lemma~\ref{zd2x-three-adic} using the entire real descent
image. In the ordering of~\eqref{zd2x-root-intervals},
put \(\alpha_i=-3\theta_i\). Then
\[
 \alpha_1<0<\alpha_3<3<\alpha_2.
\]
The real points of the monic quintic lie above
\([\alpha_1,0]\), \([\alpha_3,3]\), and
\([\alpha_2,\infty)\). Away from endpoints the signatures
of \((x,x-3;x-\alpha_1,x-\alpha_2,x-\alpha_3)\) are
\begin{equation}
 \begin{array}{c|cc|ccc}
  \text{interval}&x&x-3&x-\alpha_1&x-\alpha_2&x-\alpha_3\\ \hline
  \text{lower}&-&-&+&-&-\\
  \text{middle}&+&-&+&-&+\\
  \text{upper}&+&+&+&+&+
 \end{array}
 \label{zd2x-real-table}
\end{equation}
The lower and middle vectors generate a four-element sign
group, and their first two sign pairs are independent.
Only its identity has both first signs positive.

This restriction holds for all \(J(\mathbb R)\), not just
for the Abel--Jacobi images of real points. Move a real
divisor representative off the branch points and \(O\) as
in~\eqref{zd2x-delta-kernel}. Real closed points contribute
the vectors~\eqref{zd2x-real-table}. A nonreal closed point
is a complex conjugate pair, whose norm of \(x-r\) at any
real root \(r\) is a positive absolute-value square.
Products and inverse powers still lie in the same sign
group. This is also the explicit real-image description
of \cite[Lemma~4.8]{StollTwoDescent2001}.

Consequently an actual image \((1,1,u)\) has all three
signs of \(u\) positive. A \(K\)-square cannot change
those signs. Lemma~\ref{zd2x-field} therefore makes \(u\)
a square. We have proved
\begin{equation}
 \delta(J(\mathbb Q))=
 \langle\delta(D_0),\delta(D_3)\rangle,\qquad
 |J(\mathbb Q)/2J(\mathbb Q)|=4.
 \label{zd2x-four-global}
\end{equation}
The two known images are independent, so equality, rather
than merely an upper bound, holds.

By Mordell--Weil, \(J(\mathbb Q)\) is finitely generated.
For such a group of rank \(r\), the quotient by doubling
has size \(2^r\) times the size of its two-torsion subgroup.
That subgroup has order four here, by the rational
branch-factor calculation or Lemma~\ref{z2x-torsion-local}.
Thus~\eqref{zd2x-four-global} forces \(r=0\).
The complete torsion bound in
Theorem~\ref{z2x-finite-jacobian} then gives
\(J(\mathbb Q)\simeq(\mathbb Z/2\mathbb Z)^2\).
The rank deduction itself did not use a reduction count.

For completeness, the rational-point conclusion follows
by the explicit degree-two argument, not simply by
counting the group. The map \(R\mapsto[R-O]\) is injective.
The fourth possible group element \(D_0+D_3\) cannot be
its image: otherwise \(R+O\sim(0,0)+(3,0)\).
The canonical class is \(2O\). The functions with sole
pole at \(O\) of order at most two are the span of \(1,X\),
since the pole orders of \(X,Y\) are two and five and the
even and odd orders in \(a(X)+b(X)Y\) cannot cancel.
No such linear function vanishes at both zero and three.
Hence \((0,0)+(3,0)\not\sim2O\).
Riemann--Roch gives \(h^0((0,0)+(3,0))=1\);
its sole effective representative cannot be \(R+O\).
Thus \(C(\mathbb Q)=\{(0,0),(3,0),O\}\), which transfers
under \(s=-X/3\) to the stated three points.
\end{proof}

\begin{corollary}[The six squared-unit boundary points]
\label{zd2x-six-boundary}
The common-source curve \(D_{3,\zeta^2}\) has exactly six
rational points. They lie above the source values
\(s=0,-1,\infty\), with \(v=0\) and \(w=\pm2\) at each.
Their conic parameter is \(t=-2\) when \(w=2\) and \(t=0\)
when \(w=-2\). Its positive rational locus is empty.
\end{corollary}
\begin{proof}
Every rational point maps to the first quotient \(H\).
At each of its three rational points \(B_0=0\) and
\(A_0\ne0\) homogeneously, hence \(r=-1,v=0\).
The second square has the two values \(w=\pm2\).
It is unramified on \(H\) at these points, being a unit
square with two distinct roots. Thus each supplies
exactly two rational points on the normalization, and
there are no others. The values of \(t\) follow from
Proposition~\ref{z2x-model}. None satisfies the strict
positive condition~\eqref{z2x-positive}.
\end{proof}

\begin{theorem}[Exclusion of the actual primitive square--cube intersection]
\label{zd2x-actual-exclusion}
There are no positive integers \(a,b,U,Q\) with
\(\gcd(a,b)=1\) satisfying
\begin{equation}
 a^2+ab+b^2=U^2,\qquad
 a^4+3a^3b+5a^2b^2+3ab^3+b^4=Q^3.
 \label{zd2x-actual-system}
\end{equation}
\end{theorem}
\begin{proof}
Put \(c=a+b\), \(M=a^2+ab+b^2\), and denote the quartic
by \(F(a,b)\). The actual second Eisenstein element is
\[
 z=ab+M\zeta,\qquad N(z)=(ab)^2+abM+M^2=F(a,b).
\]
It is primitive because \(\gcd(ab,M)=1\). Moreover
\(F(a,b)\equiv(a^2+b^2)^2\equiv1\pmod3\) for a
primitive pair. Thus this element is unramified at three.

To retain the exact integral profile, use unique
factorization in \(\mathbb Z[\zeta]\), as in
Section~\ref{dc-section}. Its Euclidean norm proof is
elementary: rounding the two real coordinates of a
complex quotient in the basis \(1,\zeta\) leaves norm
at most \(3/4<1\). An inert prime in \(N(z)\) would
divide both integer coordinates. Both orientations
of a split prime in \(z\) would do the same. Hence
exactly one split orientation occurs at each norm
prime, and its exponent is the corresponding
rational exponent in \(F(a,b)\). Since this norm is
a cube, there is an exact integral equality
\begin{equation}
 ab+M\zeta=\epsilon(S+T\zeta)^3,\qquad
 \epsilon\in\{1,\zeta,\zeta^2\},
 \label{zd2x-actual-cube}
\end{equation}
with \(\gcd(S,T)=1\), an unramified root, and root
norm \(Q\). The original six units reduce to these
three by absorbing their sign into the cube root.
No nonunit residual or projective content factor has
been removed. Also \(Q>1\), since \(F(a,b)\ge13\).

The first norm is prime to three. Indeed, if
\(a\not\equiv b\pmod3\), it is a unit; otherwise
primitivity makes both units and
\(M=(a-b)^2+3ab\) has three-adic valuation one.
The latter is impossible for a square. Nor can
\(a=b\), since then the primitive seed is \((1,1)\)
and \(M=3\).
Consequently the rational numbers
\[
 r=\frac{ab}{M},\qquad v=\frac cU,\qquad
 w=\frac{a-b}{U}
\]
satisfy
\[
 0<r<\frac13,\quad v>0,\qquad
 v^2=1+r,\quad w^2=1-3r .
\]
The strict inequality follows from
\(M-3ab=(a-b)^2>0\).
Equation~\eqref{zd2x-actual-cube} supplies the very
same projective source \([S:T]\) in
\(D_{3,\epsilon}\), whose target denominator is
the nonzero integer \(M\). Both square roots are
nonzero, so this is its smooth positive chart.
Its parameter is exactly
\[
 t=\frac{v+w+2}{v-w}=\frac{a+U}{b}>1,
 \qquad v-w=\frac{2b}{U}\ne0 .
\]
This construction does not pair unrelated quotient
points or omit an infinite root source.

For \(\epsilon=1\), Section~\ref{cl-section}
excludes the entire three-adic locus. In the present
integral setting its second coordinate
\(3ST(S+T)\) also directly contradicts \(3\nmid M\).
For \(\epsilon=\zeta\), Theorem~\ref{cb-twelve-points}
gives only the twelve boundary points with \(r=0\).
For \(\epsilon=\zeta^2\),
Corollary~\ref{zd2x-six-boundary} gives only the six
boundary points with \(r=-1\).
Each contradicts the displayed strict positive
interval. These three cases exhaust all units
without a conjugacy identification, proving
\eqref{zd2x-actual-system}.
\end{proof}

\begin{corollary}[Same-coefficient propagation and its exact numerical scope]
\label{zd2x-propagation}
For every odd positive integer \(g\) divisible by
three and every \(\epsilon\in\{1,\zeta,\zeta^2\}\),
\(D_{g,\epsilon}\) has no rational point with finite
\(t>1\). The same holds for its Gaussian cover.
Numerically, for any positive integers \(Q,g\) with
\(g\) divisible by three, there is no positive primitive pair with
\(M\) a square and \(F=Q^g\).
\end{corollary}
\begin{proof}
For the geometric claim let \(k=g/3\). The degree-\(n\)
projective power map \(L_{\epsilon,n}\) is the
coordinate ratio of \(\epsilon(S+T\zeta)^n\).
Its two coordinates have no common projective zero.
The exact power identity gives
\[
 L_{\epsilon,g}=L_{\epsilon,3}\circ L_{1,k}.
\]
On the dense fiber-product chart,
\((t,s)\mapsto(t,L_{1,k}(s))\) defines a
nonconstant rational map \(D_{g,\epsilon}\to
D_{3,\epsilon}\), retaining the same \(t\).
Equivalently it retains both \(v,w\) and changes
only their common source. The odd-index curves
are geometrically connected by
Section~\ref{hg-section}; the function-field map
extends over all points of their smooth
projective models by
\cite[Lemma~53.2.2 and Theorem~53.2.6]{StacksCurves2026}.
Thus a rational point with finite \(t>1\), even
with an infinite or ramified root input, would
map to the empty positive cubic locus. The
forgetful map proves the Gaussian assertion.
No count of the remaining boundary preimages
at larger \(g\) is asserted.

For the numerical statement write
\(F=(Q^{g/3})^3\) and apply
Theorem~\ref{zd2x-actual-exclusion}. This also
allows even \(g\), without making a connectedness
claim for an even-index geometric model.
\end{proof}

\paragraph{Ordinary inputs, finite evidence, and remaining scope.}
The descent kernel and its naturality in
\eqref{zd2x-delta-kernel} are standard external
ordinary inputs from \cite{StollTwoDescent2001}.
The curve-specific field, valuation, local-image,
signature and point-set arguments are proved
above. The finite tables of
Theorem~\ref{z2x-finite-jacobian} bound full
torsion separately; they are not rank evidence.
No complete Jacobian descent or rational-point
classification is claimed formalized in Lean.

In particular \(M=R^h,F=Q^g\), with positive integers
\(R,Q,h,g\), even \(h\), and \(3\mid g\), is excluded for this
positive primitive family even when \(h,g\)
have no common prime factor. The exact total
perfect-power conditions are what matter.
The ramified first shape \(M=3R^h\) with \(h\)
even is not a square and is not covered.
A nonunit second residual is covered only
when its total \(F\) is still a cube; arbitrary
moving residuals, other indices and uniform
point-height estimates remain open.
This fixed-family exclusion is not a
classification of all primitive ABC triples.
